\paragraph{The cost of imposing the crystalline conditions}
For each \(v\mid\ell\), take the decomposition group \(\Gamma_v\)
and let \(d_v=[G_{\mathbb Q_\ell}:\Gamma_v]\), viewing the group as
an open subgroup of \(G_{\mathbb Q_\ell}\).
The local lattice class \(\mathcal C_v\) is specified by the invariance
of spanning periods over \(B_{\max}\), the set of Hodge--Tate jumps,
and polynomial annihilators of inertia. At \(v\nmid M\) the class is
crystalline with no polynomial condition on inertia. At \(v\mid M\)
it is potentially crystalline, and the inertia polynomials are those
common to all representations of the prescribed type.
Use the common jumps of all the \(N_x\).
Thus the lattices of the \(N_x\) belong to this class.
The direct-sum decomposition at \(\pi\) puts the main block and the
Tate-character lattices in the same class, whose allowed jumps come
from the actual jumps of these summands. Since the inverse twist is
trivial on \(\Gamma_v\), the main summand has the original jumps.
\leanlinks{\href{\upstream/grh/src/New/Congruences/FirstOrder/TraceRelations.lean\#L757}{The common class of local lattices}}

The crystalline cocycles \(Z_f(L)\) of a lattice \(L\) are those that
become coboundaries in \(B_{\max}\otimes L\).
Let \(\delta_v(L)\) be the normalized length of the quotient of
\(Z_f(L^\vee(1))\) by the cocycles whose cup products with the
crystalline cocycles of the lattices in \(\mathcal C_v\) are
coboundaries. The pairings are taken at finite coefficient level,
using forms equivariant modulo \(\ell^u\) whose lifted coboundary
divided by \(\ell^u\) is crystalline.

Let \(c\) be a cocycle modulo \(\ell^u\) respecting the linear
relations among \((N_x(g)-1)_x\). If \(L_0\) denotes the span of
these families, the assignment \((N_x(g)-1)_x\mapsto c(g)\) is
an equivariant linear map on \(L_0\). Equivariance follows from
\(N_x(h)(N_x(g)-1)=(N_x(hg)-1)-(N_x(h)-1)\) and the cocycle rule.
Put \(\mathcal B=\prod_xM_n(\mathcal O)\).
Artin--Rees applied to the finite \(\mathcal O\)-modules
\(L_0\subseteq\mathcal B\) and the principal ideal \((\ell^u)\)
gives \(L_0\cap \ell^{u k_0}\mathcal B\subseteq\ell^uL_0\) for
\(k_0\) sufficiently large.
Thus the map extends by zero on \(\ell^{u k_0}\mathcal B\).
Its graph injects into the first component. Projection onto all
matrix columns gives an equivariant injection into a sum of the
local lattices of the \(N_x\), so the graph is in \(\mathcal C_v\).
The graph cocycle with first component \(g\mapsto(N_x(g)-1)_x\)
becomes the coboundary of \(\ell^{u k_0}(I)_x\) after
multiplication by \(\ell^{u k_0}\). Hence it is crystalline, and
its second component is the local restriction of \(c\).

Lift the second-component map to a linear map \(q\) between lattices.
Its coboundary \(gq-q\) is divisible by \(\ell^u\), and the quotient
is a cocycle that becomes a coboundary over \(B_{\max}\).
The evaluation pairing transferred through \(q\) is therefore among
the finite-coefficient pairings above. Apply local Tate duality to
\(L/\ell^u L\) and its Tate dual over the open local decomposition
group. Using the orthogonality of the crystalline local conditions
under local duality, the order of the quotient of the cocycles just
constructed by the crystalline image divides the order of the
quotient of \(Z_f(L^\vee(1))\) by its orthogonal part.
Here evaluation over \(\mathcal O\) is composed with a generator
of its \(\mathbb Z_\ell\)-dual to give a perfect pairing over
\(\mathbb Z_\ell\). Divisibility of orders gives the length bound.
The reduction of an integral crystalline cocycle obeys all its
integral linear relations, so imposing those relations on a
finite-coefficient cocycle costs at most the same length.
The colength of imposing the conditions at all \(v\mid\ell\) is
bounded by the sum of the separate colengths.
\leanlinks{\href{\upstream/grh/src/New/Congruences/LocalLoss/Basic.lean\#L386}{The graph of a cocycle and Artin--Rees}
  \leanlinksep \href{\upstream/grh/src/New/Congruences/LocalLoss/TateDualBound.lean\#L101}{Finite-coefficient duality and the colength}}

To bound the length on the dual side, write \(n_L=\operatorname{rank}L\),
\(q_L=1\) for a main block and \(q_L=t\) for a character.
Writing \(\operatorname{def}_\ell=L_\ell/(4f)\), the required estimate is
\[
 \delta_v(L)\le d_v n_L\operatorname{def}_\ell/q_L .
\]
For \(v\mid M\) the inertia spectra provide the annihilator
\(X^{\operatorname{lcm}_{1\le i\le r}(\ell^i-1)}-1\) and a monic
polynomial over \(\mathbb Z_\ell\) of degree at most \((r+t)r\).
Since \(\ell\equiv1\) modulo the conductor, the coefficient
embeddings at the same \(v\) read the same embedding of \(E\).
The tame period calculation for the opposed jump intervals and
these annihilating polynomials sends the crystalline cocycles into
the orthogonal part on multiplication by \(\ell^{k_\mathrm{tame}}\),
where
\[
 k_\mathrm{tame}=3r^2((r+t)r+1)^2+150(\Delta^2+4t+2).
\]
By the Bloch--Kato dimension formula, the rank of the crystalline
cocycles is at most \(2d_v n_L\), giving length at most
\(2d_v n_L k_\mathrm{tame}\). The size hypotheses and \(v_\ell(M)\ge1\)
give \(2k_\mathrm{tame}\le\operatorname{def}_\ell/q_L\).

For \(v\nmid M\) write \(H=\Delta\), let \(b_{\mathrm{sm}}\) denote the smaller weight,
and let \(j\) be the lift index. Put \(Q_0=kt\), \(\eta=H^2+1\), \(W=4t\), and
\(\delta=2jH+H-2t\). The near interval is \([-\eta,\eta]\). The far intervals
have width \(W\) and start at \(b+\delta\) or \(b_{\mathrm{sm}}+\delta\),
with the signs read at the embedding. These starting points are the lift weights
minus \(2t\). The original far jumps are at least \(t\) from the endpoints, so
\(u\mapsto-u-1\) takes them into the opposite interval of the same width.
In the main block the negative jumps are at most \(-H/2\), and there are at most
\(\mu=\lfloor n_Lm/(kq_L)\rfloor\) dual jumps in the near interval.
A character has only the dual jump \(-a-1\) in \([-\eta,0]\).

The Kummer chart \(W(\mathbb F_{\ell^{f_v}})[[z]]\subset A_{\inf}\) has ramification
index \(e_v\), with \(e_v f_v=d_v\) and \(e_v\le H\).
Group the coefficient embeddings by the Teichmüller representatives of the residue-field units
and read the jumps at each unit. If \(\ell\) is prime to the cyclotomic conductor, the chart is
unramified and the representative determines the embedding of \(E\).
Otherwise \(\ell\le Q_0\) and all embeddings with the larger weight have one sign.
When \(\ell\le Q_0\), no three positions of the larger weight follow each other with both
Frobenius gaps at most \(J_0=\lfloor\log_\ell H^3\rfloor\). This is the condition on the
chosen set of large-weight places transported to the coefficient embeddings.
\leanlinks{\href{\upstream/grh/src/New/Congruences/LocalLoss/CrystallineDefect.lean\#L217}{Reading the weights and embeddings on the Kummer chart}}

Put \(\lambda=v_\ell((2b)!\Pi_b^{H^3})\) and \(K_0=\lfloor2\lambda/Q_0\rfloor\).
The factorial gives \(2b<\ell(\lambda+1)\), and division gives \(2\lambda<(K_0+1)Q_0\).

\begin{itemize}
\item If \(\ell\le Q_0\), the weight sizes and these inequalities give
\(2(b+\delta)<H^2K_0\). Include the far interval of smaller weight in the enlarged
near interval \([-2b_{\mathrm{sm}},2b_{\mathrm{sm}}]\).
Let \(d=0\) if \(e_v=1\) and \(d=v_\ell(e_v)+1\) otherwise, the chart's bound for the different.
Then
\[
 \begin{split}
 2e_v(b+\delta+10b_{\mathrm{sm}})&\le
      \ell^{J_0+1}\,2K_0,\\
 20(J_0+1+d)\,2b_{\mathrm{sm}}+
       \lfloor\log_\ell(2(b+\delta+10b_{\mathrm{sm}}))\rfloor&\le K_0 .
 \end{split}
\]
The first uses \(\ell^{J_0+1}\ge H^3\). The second uses \(J_0+1+d\le H\) and
\(41 b_{\mathrm{sm}}H\le K_0\). Take as cuts the positions followed by \(J_0\)
consecutive positions of the smaller weight. Sparsity supplies a cut at most \(2J_0+1\)
steps back. The comparison uses constant signs in the ramified case, and
\(10\ell^{J_0}\,2b_{\mathrm{sm}}\le b+\delta\) in the unramified case.
Annihilating the period cokernels at the cuts and comparing the extension points
from the cuts gives the cup-product exponent \(10^3(2K_0+(J_0+1)K_0)\le10^4K_0m\).
\leanlinks{\href{\upstream/grh/src/New/Congruences/LocalLoss/CrystallineCocycles.lean\#L333}{The exponent comparison at small primes}
  \leanlinksep \href{\upstream/grh/src/New/Congruences/Lifting/CutExponentBounds.lean\#L805}{Cuts and cup products for sparse larger weights}}

\item If \(\ell>Q_0\) and \(\lambda\ge H^3\), the chart is unramified and \(\ell\ge7\).
The two inequalities and the weight sizes give \(b+\delta<K_0\ell^2\),
\(2(b_{\mathrm{sm}}+\delta)<K_0\ell\), and \(W+4\eta\le K_0\).
Allow a pole shift of \(\max(\lfloor(b+\delta)/\ell\rfloor-7K_0,0)\) where the next position
has the larger weight, and shift zero elsewhere. A position is a cut when its shift is less
than the start of its far interval. This includes all positions of the larger weight and
all positions whose successor has the smaller weight, so there is a cut one or two steps back.
The cup-product exponent is
\[
 10^4\bigl(K_0+\lfloor\log_\ell(2(b+\delta+W+4\eta))\rfloor\bigr)
 \le10^4K_0m.
\]
In both cases the indicated power carries crystalline cocycles into the orthogonal
submodule. Their rank is at most \(2d_v n_L\), which bounds the colength after multiplying
by the exponent.
\leanlinks{\href{\upstream/grh/src/New/Congruences/LocalLoss/CrystallineCocycles.lean\#L425}{The exponent at large primes and large loss}
  \leanlinksep \href{\upstream/grh/src/New/Congruences/Lifting/CutExponentBounds.lean\#L991}{Cuts and cup products for arbitrary patterns}}

\item If \(\ell>Q_0\) and \(\lambda<H^3\), a multiple of \(\ell\) in
\((b-30H^3,b+30H^3]\) would contribute at least \(H^3\) to the valuation of \(\Pi_b^{H^3}\).
Also \(2b<\ell(\lambda+1)\) and the weight sizes imply \(\ell\ge9b_{\mathrm{sm}}\).
With \(A=2(\lambda+1)\), there are therefore no nonzero multiples of \(\ell\) within
\(6(A+\eta+W)\) of \(0\) or the starts of the far intervals with either sign.
The upper endpoints of the unsigned far intervals plus \(\eta\) are at most \((\ell-3)A\).
Under these conditions, at each residue unit take the submodule of constants with extension
points of the prescribed depth, then intersect the pullbacks over all units.
The cup product is a coboundary on this submodule and its colength is at most
\(11d_v(\mu(A+1)+n_L(W+1))\), counting length \(A+1\) for each near jump and \(W+1\)
for all jumps, summed over the coefficient embeddings.
The defect is zero for the characters, whose dual jumps lie in \([-\eta,0]\).
It is also zero for a main block when \(\lambda<\lfloor H/5\rfloor\), since the two sets
of jumps avoid \([-A,0)\). In the remaining case use \(A+1\le6\lambda\) and \(W+1\le K_0\).
\leanlinks{\href{\upstream/grh/src/New/Congruences/Lifting/FewLowJumps.lean\#L945}{The submodule of constants with extension points of the prescribed depth}
  \leanlinksep \href{\upstream/grh/src/New/Congruences/LocalLoss/CrystallineCocycles.lean\#L458}{The nondivisibility condition and the colength for small loss}}
\end{itemize}

In all three cases the
bound per unit of local degree is
\[
 10^5\bigl(n_L m\lfloor 2\lambda/(kt)\rfloor+\mu\lambda\bigr)
 \le 3\cdot10^5\,n_L m\lambda/(kq_L)
 \le n_L\operatorname{def}_\ell/q_L.
\]
The last comparisons use \(1\le q_L\le t\).
Now \(\sum_{v\mid\ell}d_v\le2f\) and over the summands
\(\sum n_L/q_L=r'+t'/t\le2r'\). Thus all the dual defects
together are at most \(r'L_\ell\).
\leanlinks{\href{\upstream/grh/src/New/Congruences/LocalLoss/CrystallineDefect.lean\#L487}{Substitution in the tame and crystalline cases}
  \leanlinksep \href{\upstream/grh/src/New/Congruences/LocalLoss/CrystallineCocycles.lean\#L571}{Three cases in the crystalline period calculation}}

A cocycle annihilating the main-block rows vanishes on inertia
where all \(N_x(g)-1\) vanish. At the Taylor--Wiles places and
their conjugates it vanishes on the new tame generators.
The main block is unramified there, so this last condition
gives vanishing on the entire inertia groups.
Imposing the crystalline conditions leaves a module contained
in the Selmer group unramified outside \(\ell,M,Q,\bar Q\)
and vanishing at \(\gamma\). With \(\ell^u\)-torsion coefficients
the length of the crystalline part is at most \(\mathrm{sel}\),
and its colength is at most the sum of the main-block defects.
The bound is uniform in \(u\), and hence also bounds the length
for divisible coefficients, the union of these torsion modules.
This accounts for \(b_{\mathrm{main}}\).
\leanlinks{\href{\upstream/grh/src/New/Congruences/FirstOrder/RibetElement.lean\#L198}{Length of the main block with divisible coefficients}}

\paragraph{Parity and length of the character cocycles}
For character cocycles, vanishing at the old inertial generators
allows ramification only above \(\ell\) and at the new places and
their conjugates. Take \(A_a\) to be the inertia groups above \(\ell\)
if \(a=0\), and the elements in the local decomposition groups
with trivial cyclotomic character if \(a>0\).
In the Bloch--Kato calculation for Tate twists, crystalline
classes for the trivial character are unramified and those for
\(\mathcal O(-a)\) with \(a>0\) are coboundaries over the fraction
field. The integral cocycles therefore vanish at the prescribed
elements, and finite-coefficient cocycles satisfying the
corresponding linear relations do so as well.
Thus the cost of vanishing on \(A_a\) is at most the crystalline
colength.

The elements \(\tau_j,\phi_j\) and their conjugates act trivially
on the finite coefficients. The cocycle rule gives
\(c(g^{-c})=-c(g^c)\). Agreement of \(c(\tau_j)\) with
\(c(\tau_j')\) is therefore equivalent to vanishing of
\(c^+=c+c\circ c_E\) at \(\tau_j\), where \(c_E(g)=g^c\),
and likewise at Frobenius elements.
The condition on even cocycles at the ordered auxiliary primes
now gives \(C(c(\tau_j)-c(\tau_j'))=0\) under the agreements at
the earlier places.

The submodule vanishing on all new inertia is unramified away
from \(\ell\). Its members vanish at \(\gamma\), so
\((\chi_a(a_{\mathrm c})-1)c(a_{\mathrm c})=0\).
If \(a=0\), the condition at \(A_0\) then gives homomorphisms
unramified everywhere. For \(a>0\), use the map
\(c\mapsto c\circ c_E-c\) to the odd classes.
In its kernel \(c\) is even. Since \(c^+(\phi_j)=0\) for
\(j\in J_a\), the reading of twice the values of the even
part at the test points implies \(4c(u)=0\) at each test point \(u\).
Put \(d_\chi=\chi_a(a_{\mathrm c})-1\). The cocycle \(4d_\chi c\)
vanishes at the test points and \(a_{\mathrm c}\), hence is zero
by detection. Since \(4d_\chi^2\mid C\), the kernel is killed
by \(C\). Thus the length modulo \(C\)-torsion is at most
the length of the odd classes. This gives \(b_a\) as the sum
of the crystalline colength and the length of the odd classes.
A cocycle that vanishes on the new inertia, test points, and
\(a_{\mathrm c}\) is zero by the same unramifiedness and detection.
This proves detection by \(s\).
\leanlinks{\href{\upstream/grh/src/New/Congruences/FirstOrder/TateTwistCocycles.lean\#L156}{The even kernel and the map to odd classes}
  \leanlinksep \href{\upstream/grh/src/New/Congruences/FirstOrder/TraceRelations.lean\#L560}{Application to the tied-row cocycles}}
